\documentclass[a4paper,12pt,oneside]{article}
\usepackage{discrete-practice}

% Источник условий: https://discrete-math-practics.yonote.ru/share/8fc01ef7-c271-4ee8-a763-121cc7eeaf6a/doc/dz4-kombinatorika-srok-do-05052026-2355-XXkPjzeVeC
\begin{document}
\HomeworkHeader{2}{4}{Комбинаторика}
\noindent Срок сдачи: 05.05.2026, 23:55.\hfill
\PracticeLink{Условия в Yonote}{https://discrete-math-practics.yonote.ru/share/8fc01ef7-c271-4ee8-a763-121cc7eeaf6a/doc/dz4-kombinatorika-srok-do-05052026-2355-XXkPjzeVeC}

\medskip\noindent\textbf{Важно:} в каждом задании вкратце поясните, почему вы выбрали тот или иной комбинаторный объект для вычисления (важен ли порядок, есть повторения или нет и т.\,п.) / по какому принципу вы вели расчёты.

\Task{Задание 1}
Сколько различных слов можно получить, переставляя буквы слова «инвариант» так, чтобы каждая буква «а» шла непосредственно после некоторой «н»?

\Task{Задание 2}
17 конфет необходимо распределить по трем ящикам. Сколькими способами это можно сделать при условии, что все конфеты одинаковые?

\Task{Задание 3}
Сколькими способами могут быть размещены дни рождения 12 человек в году, считая, что в нем 365 дней?

Считаем, что «порядок» важен. То есть вариант 1: у человека A день рождения 01.02, у человека B -- 01.03 и вариант 2: у человека A день рождения 01.03, у человека B -- 01.02 -- это 2 разных варианта.
\begin{alphlist}
  \item Во скольких случаях все дни рождения попадут на разные дни года?
  \item А во скольких на разные месяцы? (При этом важно, в какой именно день месяца будет день рождения.)
\end{alphlist}

\Task{Задание 4}
Сколько существует положительных целых чисел, содержащих не более пяти цифр, в которых
\begin{alphlist}
  \item Первой цифрой является 3?
  \item Последней цифрой является 5?
  \item Первой цифрой является 3 или последней цифрой является 5?
  \item Ни первая цифра не равна 3, ни последняя цифра не равна 5?
\end{alphlist}

\Task{Задание 5}
Сколько трёхзначных чисел можно образовать из цифр 2, 3, 4, 5, 6, 8, 9? (В этом задании цифры могут повторяться.)
\begin{alphlist}
  \item Сколько из таких трёхзначных чисел меньше 450?
  \item Сколько среди них чётных?
  \item Сколько из них делятся на 4?
\end{alphlist}

\newpage
\Task{Задание 6}
Пусть $A=\{a,b,c,d,e,f,g,h\}$. Сколько существует
\begin{alphlist}
  \item трёхэлементных подмножеств множества $A$?
  \item пятиэлементных подмножеств множества $A$, содержащих $b$?
  \item пятиэлементных подмножеств множества $A$, не содержащих $b$?
  \item пятиэлементных подмножеств множества $A$, содержащих $c$, но не содержащих $d$ и $e$?
  \item подмножеств множества $A$, содержащих хотя бы 3 элемента?
\end{alphlist}

\Task{Задание 7}
\begin{alphlist}
  \item Дана перестановка 416532. Определите 3 предыдущих и 3 следующих перестановки.
  \item Пусть $S=\{1,2,3,4,5,6\}$. Определите 3 предыдущих и 3 следующих сочетания для сочетания $\{1,6\}$. Сочетания могут быть любой длины.
  \item Пусть $S=\{1,2,3,4,5,6\}$. Определите 3 предыдущих и 3 следующих сочетания длины 4 для сочетания $\{1,3,4,5\}$.
  % Длина 3 и четыре элемента сохранены буквально по источнику.
  \item Пусть $S=\{a,b,c,d,e,f\}$. Определите 3 предыдущих и 3 следующих сочетания с повторением длины 3 для $\{a,c,d,e\}$.
\end{alphlist}

\Task{Задание 8}
Сколько существует решений уравнения $x+y+z+k=23$, если $x\ge2$, $y\ge3$, $z\ge4$, $k\ge5$?

\Task{Задание 9}
Какой коэффициент будет стоять перед $x^8y^5$ в разложении $(3x-4y)^{13}$?

\Task{Задание 10}
Воспользуйтесь формулой Паскаля и докажите равенство:
\[
  C_{n+2}^{k+2}=C_n^k+2C_n^{k+1}+C_n^{k+2}.
\]
\end{document}
